%source mathadata
\begin{multicols}{2}

		\begin{MADenumerate}
			\item Déterminer par lecture graphique l'équation réduite de la droite $(AB)$ ci-contre.
			\item On procède maintenant comme une IA : sans regarder le graphique, on veut calculer si le point $C(2;4)$ appartient à la droite.
			\begin{enumerate}
				\item Déterminer les coordonnées du point $E$, qui a la même abscisse que $C$ et appartient à la droite $(AB)$.
				\item Comparer l'ordonnée de $C$ avec l'ordonnée de $E$.
				\item Déduire si $C$ est au-dessus ou en dessous de la droite $(AB)$.
			\end{enumerate}

			\item (Bonus) Proposer une méthode générale pour déterminer si un point est au-dessus ou en dessous d'une droite, en fonction de son équation réduite et des coordonnées du point.

		\end{MADenumerate}
		\begin{center}
			\begin{figure}[H]\centering
				\definecolor{wrwrwr}{rgb}{0.380,0.380,0.380}
				\definecolor{rvwvcq}{rgb}{0.082,0.396,0.753}
				\begin{tikzpicture}[line cap=round,line join=round,>=triangle 45,x=1cm,y=1cm,scale=0.8]
					\begin{axis}[
							x=1cm,y=1cm,
							axis lines=middle,
							ymajorgrids=true,
							xmajorgrids=true,
							xmin=-1.5,
							xmax=7.3,
							ymin=-1.4,
							ymax=7.2,
							xtick={-1,-3,...,7},
							ytick={-1,0,...,7},]
						\clip(-1.2,-1.4) rectangle (7.2,7.2);
						\draw [line width=2pt,domain=-4.8:9.3] plot(\x,{(--12-3*\x)/4});
						\draw [line width=2pt,dash pattern=on 1pt off 1pt] (2,4)-- (2,1.6);
						\begin{scriptsize}
							\draw [fill=blue!30] (0,3) circle (2.5pt);
							\draw [fill=blue!30] (0.3,3.3) node {$A$};
							\draw [fill=blue!30] (4,0) circle (2.5pt);
							\draw [fill=blue!30] (4.3,0.3) node {$B$};
							\draw [fill=blue!30] (2,4) circle (2.5pt);
							\draw [fill=blue!30] (2.3,4.3) node {$C$};
							\draw [fill=wrwrwr] (2,-0.75*2+3) circle (2.5pt);
							\draw [fill=black!30] (2.3,-0.75*2+3.3) node {$E$};
						\end{scriptsize}
					\end{axis}
				\end{tikzpicture}
			\end{figure}
		\end{center}
	\end{multicols}
